\section{Legal Value of the Optimality Certificate}
\label{sec:legal}

\subsection{The Certificate Statement}

When the \ilp{} solver returns \texttt{Optimal} with gap $= 0.0$, the
audit JSON contains:

\begin{verbatim}
"solver_status": "optimal",
"optimal_ec": 847,
"achieved_gap": 0.0,
"solver_used": "glpk",
"solver_version": "GLPK 5.0"
\end{verbatim}

This record supports the following legal assertion, which is unique among
algorithmic redistricting tools:

\begin{quote}
\textbf{Certificate of Compactness Optimality.}
\textit{The redistricting plan submitted herein is the globally optimal
solution to the edge-cut minimisation problem for the State of New Hampshire,
2020 Census, with $k = 2$ congressional districts and a population balance
tolerance of $\pm 0.5\%$.
No valid redistricting plan --- with the same census data, the same number
of districts, and the same population balance constraint --- can achieve a
lower total edge cut.
The optimality of this plan has been verified by Integer Linear Programming
with a certified zero optimality gap.
The complete ILP formulation is retained at \texttt{runs/nh\_optimal/2020/ilp/}
and can be independently verified by any LP/MIP solver.}
\end{quote}

This statement has no heuristic equivalent.
A practitioner submitting a \metis{}-based plan can claim only that it is
``among the best plans found in a search of 10{,}000 seeds'' --- a
statistical argument subject to challenge.
The ILP certificate is a mathematical proof.

\textbf{Important:} this certificate applies only when
\texttt{achieved\_gap = 0.0} (\texttt{solver\_status: Optimal}).
A \texttt{GapReached} result (e.g., gap $\leq 1\%$) supports only a
bounded-suboptimality claim: ``no valid plan has edge cut lower than the
submitted plan's EC $-$ achieved\_gap $\times$ EC''.
Practitioners must verify \texttt{solver\_status} before citing the
zero-gap certificate.

\subsection{Scope and Limitations of the Certificate}

The ILP certificate is precisely scoped and practitioners must represent
it accurately:

\paragraph{Edge cut only.}
The certificate covers edge-cut compactness (the number of cut adjacency
edges, which proxies for the total boundary length between districts and
non-adjacent counties).
It does not certify optimality for Polsby-Popper, Reock, convex-hull ratio,
or any other compactness metric.
A plan optimal by edge cut may not be optimal by Polsby-Popper; these are
different objectives.

\paragraph{Contiguity model.}
The certificate assumes MTZ contiguity constraints (connected districts).
This matches the standard redistricting contiguity requirement.
However, some jurisdictions require contiguity under a stricter definition
(e.g., no peninsulas narrower than a minimum width, or contiguity along
roads rather than adjacency edges).
Practitioners should verify that the MTZ contiguity model matches the
applicable legal standard.

\paragraph{Population balance tolerance.}
The certificate is conditional on the specified $\varepsilon_{\mathrm{bal}}$.
A certificate at $\varepsilon_{\mathrm{bal}} = 0.5\%$ does not certify
optimality at $\varepsilon_{\mathrm{bal}} = 1.0\%$; the feasible region
changes and the optimal plan may differ.
The tolerance must be disclosed alongside the certificate.

\paragraph{Census data vintage.}
The certificate is for the specified census year.
A plan optimal under 2020 Census data may not be optimal under 2010 Census
data or any future Census, because tract populations change.
The certificate is inherently vintage-specific.

\paragraph{Legal caveat.}
The ILP certificate addresses a purely mathematical question (minimum
edge cut under population balance and contiguity constraints) and does not
constitute legal advice.
Redistricting compliance requires assessment of criteria beyond edge-cut
compactness, including Voting Rights Act compliance, preservation of
political subdivisions, protection of communities of interest, and
applicable state constitutional standards, all of which are legal questions
outside the scope of this tool.
Practitioners should consult legal counsel before representing the ILP
certificate in litigation or regulatory proceedings.

\paragraph{Compactness versus partisan gerrymandering.}
Edge-cut optimality and partisan neutrality are orthogonal properties.
The ILP certificate asserts that no valid plan achieves lower edge cuts;
it does not certify that the plan is not a partisan gerrymander.
For 2-district states (NH $k=2$, RI $k=2$), the binary partition structure
means that the most compact plan by edge cut may systematically favour one
party's geography.
Practitioners should report partisan outcomes alongside the compactness
certificate and consult the full G.14 decision framework for multi-metric
evaluation.

\subsection{Independent Auditor Verification}

The retained MPS file enables a fully independent verification path:

\begin{enumerate}
\item Download the MPS file from \texttt{runs/\{label\}/\{year\}/ilp/\{node\}.mps}.
\item Load it into any LP/MIP solver (Gurobi, CPLEX, \highs{}, GLPK, SCIP).
\item Solve with optimality gap 0.
\item Confirm that the solver returns the same optimal objective value
  as reported in the audit JSON.
\item Confirm that the solver's optimal assignment matches the committed
  district plan.
\end{enumerate}

This verification requires no access to the platform source code,
no Rust toolchain, and no proprietary infrastructure.
Any solver that can read the standard MPS format can verify the certificate.

\textbf{Note (Phase~1 limitation):} MPS files are currently stored as loose
files under \texttt{runs/\{label\}/\{year\}/ilp/}.
They are \emph{not} included in the \texttt{label-verify} SHA audit chain
in Phase~1.
This means that MPS files can be modified after generation without detection
by the platform's integrity checks.
Phase~2 will add MPS file hashes to the \texttt{index.json} manifest,
enabling tamper-evident chain-of-custody.
Until then, practitioners submitting ILP certificates in legal proceedings
should independently hash the MPS file at time of generation and include
the hash in any expert report.

\subsection{Comparison with Ensemble-Based Expert Reports}

Expert witness reports in redistricting litigation commonly use
ensemble analysis: generate 10{,}000--100{,}000 random valid plans
and report where the disputed plan falls in the distribution.
This approach is well-established (Mattingly and Vaughn, DeFord et al.)
but has a structural limitation: an ensemble can show that a plan is
``unusually partisan'' but cannot prove that a more compact plan exists
that achieves the same partisan outcome.

The ILP certificate addresses a complementary question: it proves that
no more compact plan exists under the edge-cut metric.
In a litigation context, ILP and ensemble analysis serve different
arguments:

\begin{itemize}
\item Ensemble analysis: ``This plan's partisan outcome is in the 1st
  percentile of 50{,}000 valid plans'' (statistical).
\item ILP certificate: ``This plan's edge-cut compactness is globally
  optimal; no valid plan is more compact by this metric'' (exact).
\end{itemize}

Used together, these provide a stronger legal argument than either alone.

The ILP certificate certifies edge-cut compactness optimality under the
adjacency-graph model.
This is a necessary but not sufficient condition for legal compactness:
courts in most jurisdictions apply Polsby-Popper, Reock, or convex-hull
compactness --- none of which equals edge cut.
The plan may still be challenged on other compactness grounds.
Expert witnesses should qualify the certificate's scope precisely:
``this plan is provably optimal among all valid plans by the edge-cut
measure; compliance with Polsby-Popper or other legal compactness
standards requires separate analysis.''
